\section{Conclusion}
\label{sec:conclusion}

We presented a systematic study of how signal dilution affects LLM steganographic capabilities, testing three frontier models across six encoding schemes in a unified framework covering encoding, informed extraction, and blind detection. Our central finding is a cliff effect: all three models achieve 92--100\% per-character accuracy on acrostic encoding but collapse to 0--43\% on every diluted scheme, with 0\% exact match across all diluted conditions. This transition is abrupt rather than gradual, suggesting that acrostic steganography is a memorized pattern rather than an instance of a general compositional steganographic capability. Blind detection fails entirely at all dilution levels, with each model applying a fixed response bias regardless of whether hidden messages are present.

These results carry a clear practical message: current LLMs can neither produce nor detect steganographic content beyond the well-practiced acrostic pattern. For AI safety, this means the immediate risk of covert LLM communication through diluted encoding is low, but so is the reliability of LLM-based steganography monitors.

Several directions remain open. Chain-of-thought prompting or explicit counting scaffolds may close part of the capability gap on diluted schemes. Fine-tuning on diluted encoding tasks would establish whether the cliff reflects a fundamental architectural limitation or a training data gap. Tool-augmented LLMs with access to word counters could bypass the counting bottleneck entirely. Finally, replacing binary detection with forced-choice or information-theoretic approaches may yield more calibrated detection systems. Understanding where the cliff can and cannot be moved is essential for anticipating the trajectory of steganographic capabilities as models continue to advance.
